% Part: first-order-logic
% Chapter: axiomatic-deduction
% Section: proof-theoretic-notions

\documentclass[../../../include/open-logic-section]{subfiles}

\begin{document}

\iftag{FOL}
      {\olfileid[pt]{fol}{axd}{ptn}}
      {\olfileid[pt]{pl}{axd}{ptn}}

\olsection{Noções da teoria da prova}

\begin{explain}
Assim como definimos várias noções semânticas importantes
(\iftag{FOL}{validade}{tautologia}, consequência, satisfatibilidade),
agora definimos as \emph{noções da teoria da prova} correspondentes.
Elas não são definidas por meio da satisfação de !!{sentence}s em
!!{structure}s, mas por meio da !!{derivability} ou da
!!{nonderivability} de certas fórmulas. A descoberta de que essas
noções coincidem foi importante. Essa coincidência é o conteúdo dos
teoremas da \emph{correção} e da \emph{completude}.
\end{explain}

\begin{defn}[!!^{derivability}]
!!^a{formula}~$!A$ é \emph{!!{derivable}} a partir de $\Gamma$, o que
escrevemos $\Gamma \Proves !A$, se existe !!a{derivation} a partir
de~$\Gamma$ que termina em~$!A$.
\end{defn}

\begin{defn}[Teoremas]
!!^a{formula}~$!A$ é um \emph{teorema} se existe !!a{derivation} de
$!A$ a partir do conjunto vazio. Escrevemos $\Proves !A$ se $!A$ é um
teorema e $\Proves/ !A$ se não é.
\end{defn}

\begin{defn}[Consistência]
Um conjunto $\Gamma$ de !!{formula}s é \emph{consistente} se, e somente
se, $\Gamma\Proves/ \lfalse$; caso contrário, é \emph{inconsistente}.
\end{defn}

\begin{prop}[Reflexividade]
\ollabel{prop:reflexivity}
Se $!A \in \Gamma$, então $\Gamma \Proves !A$.
\end{prop}

\begin{proof}
  A !!{formula}~$!A$ por si só é !!a{derivation} de~$!A$ a partir de~$\Gamma$.
\end{proof}

\begin{prop}[Monotonicidade]
\ollabel{prop:monotonicity}
Se $\Gamma \subseteq \Delta$ e $\Gamma \Proves !A$, então $\Delta
\Proves !A$.
\end{prop}

\begin{proof}
Toda !!{derivation} de $!A$ a partir de $\Gamma$ também é !!a{derivation}
de $!A$ a partir de~$\Delta$.
\end{proof}

\begin{prop}[Transitividade]
\ollabel{prop:transitivity}
Se $\Gamma \Proves !A$ e $\{!A\} \cup \Delta \Proves
!B$, então $\Gamma \cup \Delta \Proves !B$.
\end{prop}

\begin{proof}
  Suponha $\{!A\} \cup \Delta \Proves !B$. Então existe
  !!a{derivation} $!B_1$, \dots, $!B_l = !B$ a partir de~$\{!A\} \cup
  \Delta$. Alguns passos dessa !!{derivation} serão corretos por uma
  regra que se refere a uma linha anterior~$!B_i = !A$. Por hipótese,
  existe !!a{derivation} de~$!A$ a partir de~$\Gamma$, isto é,
  !!a{derivation}~$!A_1$, \dots, $!A_k = !A$ na qual cada $!A_i$ é um
  axioma, !!a{element} de~$\Gamma$ ou é correto por uma regra de
  inferência. Considere agora a sequência
  \[
  !A_1, \dots, !A_k = !A, !B_1, \dots, !B_l = !B.
  \]
  Esta é uma !!{derivation} correta de~$!B$ a partir de $\Gamma \cup
  \Delta$, pois todo $B_i = !A$ agora é justificado pela mesma regra
  que justifica~$!A_k = !A$.
\end{proof}

Observe que isso significa, em particular, que, se $\Gamma \Proves !A$
e $!A \Proves !B$, então $\Gamma \Proves !B$. Segue também que, se
$!A_1, \dots, !A_n \Proves !B$ e $\Gamma \Proves !A_i$ para cada~$i$,
então $\Gamma \Proves !B$.

\begin{prop}
\ollabel{prop:incons}
$\Gamma$ é inconsistente se, e somente se, $\Gamma \Proves !A$ para todo~$!A$.
\end{prop}

\begin{proof}
Exercício.
\end{proof}

\tagprob{FOL}
\begin{prob}
Demonstre \olref[fol][axd][ptn]{prop:incons}.
\end{prob}
\tagendprob

\tagprob{notFOL}
\begin{prob}
Demonstre \olref[pl][axd][ptn]{prop:incons}.
\end{prob}
\tagendprob

\begin{prop}[Compacidade]
\ollabel{prop:proves-compact}
  \begin{enumerate}
  \item Se $\Gamma \Proves !A$, então existe um subconjunto finito $\Gamma_0
    \subseteq \Gamma$ tal que $\Gamma_0 \Proves !A$.
  \item Se todo subconjunto finito de~$\Gamma$ é
    consistente, então $\Gamma$ é consistente.
  \end{enumerate}
\end{prop}

\begin{proof}
  \begin{enumerate}
    \item Se $\Gamma \Proves !A$, então existe uma sequência finita de
      !!{formula}s $!A_1$, \dots,~$!A_n$ tal que $!A \ident !A_n$ e
      cada $!A_i$ é um axioma lógico, !!a{element} de~$\Gamma$ ou
      decorre de !!{formula}s anteriores por modus ponens. Tome
      $\Gamma_0$ como o conjunto dos $!A_i$ que pertencem a~$\Gamma$.
      Então a !!{derivation} também é !!a{derivation} a partir
      de~$\Gamma_0$, e portanto $\Gamma_0 \Proves !A$.
    \item Esta é a contrapositiva de~(1) para o caso especial $!A
      \ident \lfalse$.
\end{enumerate}
\end{proof}

\end{document}
